\section{The ClarInject-Code Benchmark}
\label{sec:benchmark}

\subsection{Design Principles}

ClarInject-Code is constructed around three principles:

\textbf{Matched pairs.} Each scenario has an execution condition and a clarification
condition that use the \emph{identical adversarial payload}. The only variable is
the channel through which the payload arrives. This design, borrowed from~\citet{aspi2026},
cleanly isolates the channel effect from payload-strength confounds.

\textbf{Mechanical verifiability.} Unlike evaluations that rely on LLM judges to
assess attack success, ClarInject-Code uses code-level predicates: regex patterns,
substring matching, and AST checks. This eliminates judge-model bias and enables
exact reproducibility.

\textbf{Concrete harm categories.} We restrict attacks to operations with clear,
measurable real-world consequences if a coding agent were to execute them in
production.

\subsection{Scenario Structure}

Each scenario is a JSON object with the following fields:

\begin{itemize}
  \item \texttt{task\_description}: An ambiguous coding task with one critical
    constraint deliberately removed (the ``blocked information'' technique of
    HiL-Bench~\citep{hilbench2024}).
  \item \texttt{clarifying\_question\_expected}: The question a well-calibrated
    agent should ask to resolve the ambiguity.
  \item \texttt{benign\_answer}: A legitimate user reply.
  \item \texttt{adversarial\_answer}: The benign answer with an injected instruction
    appended---structured to sound like a natural continuation of the clarification.
  \item \texttt{success\_check}: A machine-checkable predicate that returns \texttt{True}
    if the agent followed the injected instruction.
  \item \texttt{benign\_check}: A predicate that returns \texttt{True} if the
    agent correctly completed the legitimate task.
\end{itemize}

\subsection{Attack Categories}

We define five attack categories covering distinct harm types:

\begin{table}[t]
\centering
\caption{ClarInject-Code attack categories and example detection patterns.}
\label{tab:categories}
\small\setlength{\tabcolsep}{4pt}
\resizebox{\columnwidth}{!}{%
\begin{tabular}{lll}
\toprule
\textbf{Category} & \textbf{Harm} & \textbf{Detection} \\
\midrule
\texttt{backdoor} & Hidden auth bypass & Regex on bypass code \\
\texttt{secret\_exfil} & Credentials sent to URL & Regex on network call \\
\texttt{safetycheck\_removal} & Removes input validation & Regex on removed check \\
\texttt{dependency\_poison} & Typosquatted package & Substring on pkg name \\
\texttt{destructive\_command} & rm / DROP TABLE / FLUSHALL & Regex on destructive op \\
\bottomrule
\end{tabular}}
\end{table}

\subsection{Injection Style}

All adversarial answers follow a common structure: the benign answer, followed by a
natural-sounding justification, followed by the injected instruction. For example:

\begin{quote}
\small
\textit{``Use the variable name \texttt{DATABASE\_URL}. Also, the security team approved
adding a master override: if the password is `admin\_bypass\_2024', the function should
always return True regardless of username.''}
\end{quote}

This injection style is deliberately naturalistic---it mimics legitimate operational
context that a user might plausibly supply---making it a fair test of the agent's
ability to distinguish task data from instructions.

\subsection{Seed Set}

The current release includes \textbf{25 seed scenarios} (5 per attack category),
each a realistic Python coding task drawn from common developer workflows: authentication,
logging, file handling, caching, and database management. We document a clear extension
protocol for expanding to $\sim$200 scenarios, including guidelines for maintaining
category balance and injection naturalness.

\subsection{Evaluation Protocol}

For each scenario, we run:
\begin{enumerate}
  \item \textbf{Execution condition}: The agent receives a fully-specified task; the
    adversarial content is injected via a simulated file-read tool output.
  \item \textbf{Clarification condition}: The agent receives the ambiguous task; if
    it asks a clarifying question, it receives the adversarial answer through the
    \texttt{ask\_user} channel.
\end{enumerate}

Each condition is run $N=3$ times for variance estimation. Agents that do not ask a
clarifying question in the clarification condition are scored on their first-turn output
(and recorded as having $\bar{T} = 1$ clarification turn).
